% =============================================================================
%  PRICAI 2026 paper draft.
%
%  PRICAI proceedings are published in Springer LNAI/LNCS, so this draft uses the
%  Springer LNCS document class. The class file is NOT redistributed here.
%
%  >>> BEFORE COMPILING <<<
%  Download the official Springer LNCS authoring kit and place the following
%  files next to this main.tex:
%      - llncs.cls
%      - splncs04.bst
%  Source: https://www.springer.com/gp/computer-science/lncs/conference-proceedings-guidelines
%  (Overleaf also ships a "Springer LNCS" template that bundles both files.)
%
%  If you do not have llncs.cls available and just want to preview the content,
%  comment out the \documentclass line below and uncomment the article-based
%  fallback block. The fallback is best-effort; the camera-ready MUST use llncs.
% =============================================================================

\documentclass[runningheads]{llncs}

% ----------------------------------------------------------------------------
% FALLBACK (uncomment this block AND comment out the \documentclass above if
% llncs.cls is unavailable). Remove before submission.
% ----------------------------------------------------------------------------
% \documentclass[11pt]{article}
% \usepackage[margin=1in]{geometry}
% \newcommand{\keywords}[1]{\par\noindent\textbf{Keywords:} #1}
% \usepackage{abstract}
% ----------------------------------------------------------------------------

% ---------------------------------------------------------------------------
% This paper contains Chinese (CJK) strings (task queries, place names).
% RECOMMENDED build engine: XeLaTeX or LuaLaTeX (handles UTF-8 CJK natively).
% A pdfLaTeX fallback via CJKutf8 is provided, but XeLaTeX is preferred.
% All CJK text is wrapped in \cn{...} so the engine switch is centralized.
% ---------------------------------------------------------------------------
\usepackage{iftex}
\ifXeTeX
  \usepackage{fontspec}
  \usepackage{xeCJK}
  % TODO: set any installed CJK font. SimSun ships with Windows; on Linux/Overleaf
  % use e.g. "Noto Sans CJK SC". Comment out if your xeCJK auto-selects a font.
  \setCJKmainfont{SimSun}
  \newcommand{\cn}[1]{#1}
\else\ifLuaTeX
  \usepackage{fontspec}
  \usepackage{luatexja-fontspec}
  \newcommand{\cn}[1]{#1}
\else
  % pdfLaTeX fallback (requires the CJK package fonts, e.g. arphic/gbsn).
  \usepackage[T1]{fontenc}
  \usepackage[utf8]{inputenc}
  \usepackage{CJKutf8}
  \newcommand{\cn}[1]{\begin{CJK}{UTF8}{gbsn}#1\end{CJK}}
\fi\fi
\usepackage{graphicx}
\usepackage{booktabs}
\usepackage{amsmath}
\usepackage{amssymb}
\usepackage{array}
\usepackage{multirow}
\usepackage{tabularx}
\usepackage{url}
\usepackage{xurl}
\usepackage{xcolor}
\usepackage[hidelinks]{hyperref}

% Keep literal model names, log labels, paths, and API identifiers intact while
% allowing TeX to wrap them at punctuation instead of protruding past the margin.
\DeclareRobustCommand{\code}[1]{\allowbreak\nolinkurl{#1}\allowbreak}
\newcolumntype{Y}{>{\raggedright\arraybackslash}X}
% Only assists ordinary prose in an exceptionally tight line; tables are fixed
% structurally below and are not scaled or clipped.
\setlength{\emergencystretch}{1.2em}

% Robust figure placeholder: if figures/figN.pdf does not yet exist, draw a
% labelled box instead so the document still compiles. Replace the PDFs to get
% the real figures. (Remove this helper once all figures are generated.)
\newcommand{\figplaceholder}[2]{%
  \IfFileExists{#1}{\includegraphics[width=#2\linewidth]{#1}}{%
    \fbox{\parbox[c][3.2cm][c]{#2\linewidth}{\centering\ttfamily\small
      [PLACEHOLDER] #1\\ generate this figure (see caption + spec)}}%
  }%
}

\setcounter{topnumber}{3}
\setcounter{bottomnumber}{2}
\renewcommand{\topfraction}{0.85}
\renewcommand{\bottomfraction}{0.70}
\renewcommand{\textfraction}{0.15}
\renewcommand{\floatpagefraction}{0.80}

\begin{document}


\title{How Often Should a Navigator Replan? An Evidence-Audited Cadence Study for a CDP Browser Agent}

% PRICAI 2026 uses single-blind reviewing — author names are NOT hidden.
\author{Youhao Zhang \inst{1} \and Zhihan Yang \inst{2}}
\institute{
Wenzhou-Kean University, College of Science and Technology, Wenzhou, China \\
\email{zhanyouh@kean.edu}
\and
Shenzhen University, Shenzhen, China \\
\email{grace.yang.xtbb@outlook.com}
}

\authorrunning{Y. Zhang and Z. Yang}
\titlerunning{Navigator Replanning Frequency for Browser Agents}

\maketitle

\begin{abstract}
Large language model (LLM) browser agents are often augmented with external
planners, but it remains unclear how often such planners should intervene during
execution. We study navigator invocation cadence in a text-only CDP browser
agent with a fixed Doubao executor. Six configurations are evaluated: no
navigator, one-shot planning, fixed-interval replanning every 1/3/5 steps, and
event-triggered replanning. Across four live web tasks with 10 repetitions per
configuration (240 runs), strict success is audited against frozen task-specific
evidence rather than agent self-report. One-shot planning and event-triggered
replanning both achieved 40/40 strict successes, whereas frequent
fixed-interval replanning reduced success on selected tasks. In this four-task
case study, event-triggered replanning had the lowest median duration on three
tasks (R-5 was slightly faster on HuggingFace). Hospital lookup provides the
clearest diagnostic case: one-shot and adaptive guidance substantially reduced
stall burden relative to the executor-only baseline. Runs were executed
sequentially on live websites with randomized condition order within each task
block; the navigator and executor use different model families. Under this
controlled deployment, the sweep yields actionable cadence guidance---one-shot
or event-triggered replanning for reliability, and conservative fixed intervals
when replanning is required---while broader model and task coverage remains
future work.
\keywords{LLM browser agents \and planner--executor \and web automation \and
agent evaluation \and CDP.}
\end{abstract}

% =============================================================================
\section{Introduction}

LLM-driven browser agents repeatedly perceive a page, select an action, and
continue until a task is complete. In deployed live-web settings, an external
navigator or planner is often added to guide a fixed executor. However, it
remains unclear how often this guidance should be refreshed during execution.
A final \texttt{done} message can also conceal missing fields, incomplete
evidence, or an incorrect recovery path, making it difficult to distinguish
apparent completion from evidence-supported task success.

We study a narrow intervention: a CDP-based browser-use agent with one fixed
executor and an optional navigator. The navigator contributes either an opening
plan, periodic tactical subgoals, or an event-triggered recovery subgoal; it never
emits browser actions. This design holds the executor, task cards, runtime, and
evidence audit fixed while varying only navigator invocation policy.

\paragraph{Research question.} Holding the executor and runtime fixed, how do
one-shot, fixed-interval, and event-triggered navigator policies change
evidence-audited completion and execution behavior on live web tasks?

\paragraph{Contributions.}
\begin{itemize}
  \item An evidence-audited cadence sweep for a fixed CDP browser-agent
  executor: E (no navigator), I (one-shot plan), R-1/R-3/R-5
  (fixed-interval replanning), and R-A (event-triggered replanning), with
  240 runs in total.
  \item A strict-success protocol that verifies task-specific evidence rather
  than agent self-report, supplemented by compact diagnostic measures of
  execution progress.
  \item Descriptive, deployment-specific evidence---interpreted as a four-task
  case study rather than a benchmark claim---that one-shot planning and
  event-triggered replanning both achieve 40/40 strict successes, whereas
  frequent fixed-interval replanning can reduce success on selected tasks.
\end{itemize}

Because the navigator (DeepSeek) and executor (Doubao) use different model
families, our estimates characterize a deployed system configuration rather
than an isolated planning effect; a same-model navigator control is left to
future work.

\section{Related Work}

\paragraph{Web agents and evaluation environments.}
Recent work has developed both benchmark environments and end-to-end systems for
LLM-driven web interaction. WebArena~\cite{zhou2023webarena} provides
self-hosted replicas of realistic websites for reproducible functional evaluation,
whereas Mind2Web~\cite{deng2023mind2web} formulates generalist web-action
prediction from demonstrations. WebVoyager~\cite{he2024webvoyager} evaluates
end-to-end multimodal agents on live websites, thereby exposing the instability
of production-web interaction. SeeAct~\cite{zheng2024seeact} emphasizes visual
grounding for action selection, while Agent-E~\cite{emergence2024agente}
combines hierarchical coordination, DOM distillation, and change observation for
web navigation. Our study is complementary: it does not propose a new general
agent architecture or benchmark, but holds one deployed CDP executor fixed and
isolates the cadence of optional navigator intervention.

\paragraph{Planning, execution, and self-correction.}
The broader agent literature has studied how reasoning can guide tool use and
recovery. ReAct~\cite{yao2023react} interleaves reasoning and actions;
Reflexion~\cite{shinn2023reflexion} uses verbal feedback for iterative
self-improvement; and Plan-and-Solve~\cite{wang2023planandsolve} separates
problem decomposition from downstream execution. Most closely related,
Plan-and-Act~\cite{erdogan2025planandact} explicitly separates a planner from
an executor for long-horizon web tasks and evaluates the benefit of stronger plan
generation. These works motivate planner--executor decomposition, but they do
not isolate how often a lightweight external navigator should be invoked while
the executor, runtime, task cards, and scoring rule remain unchanged. We study
that narrower operational question through one-shot, periodic, and
stall/revisit-triggered policies.

\paragraph{Practical web-automation systems and scalable simulation.}
Systems work also highlights deployment-oriented concerns. Steward~\cite{wang2024steward}
studies natural-language web automation on live websites, including reactive
planning/execution and runtime efficiency. Cybernaut~\cite{park2025cybernaut}
targets reliable enterprise web automation through SOP generation, DOM element
recognition, and execution-consistency assessment. WebWorld~\cite{yao2025webworld}
addresses a complementary scaling problem by learning a web world model for
synthetic trajectory generation and simulation. In contrast, our experiments use
live public sites and a fixed browser-agent stack to characterize the behavioral
trade-off introduced by navigator cadence rather than to scale training or
construct a new automation architecture.

\paragraph{Evidence-based agent evaluation.}
Reliable evaluation remains an open problem. AgentBench~\cite{liu2024agentbench}
provides a multi-environment evaluation suite, and LLM-as-a-judge
methods~\cite{zheng2023judging} offer a flexible alternative where deterministic
checkers are unavailable. Our protocol instead uses frozen task-specific hard
criteria, saved final evidence, and visited URLs where appropriate. This makes
strict success intentionally stricter than an agent-declared \code{done} signal
and supports diagnosing completion failures on live tasks without treating a
single trajectory as uniquely correct.


% =============================================================================
\section{System and Experimental Protocol}
% =============================================================================

\subsection{Agent architecture and interventions}

% IMPORTANT: replace the figure asset before submission so that its internal labels read
% ``Experimental Conditions (6)'' and ``4 Tasks $\times$ 6 Conditions $\times$ 10 Repeats = 240 Runs''.
\begin{figure}[t]
\centering
\figplaceholder{figures/fig1.png}{0.94}
\caption{System architecture. All conditions share a fixed executor LLM driving
a CDP browser loop. A navigator can provide an opening plan (I), periodic
guidance (R-$k$), or event-triggered recovery guidance (R-A) by injecting a
\code{<current\_step\_focus>} subgoal into the executor prompt.}
\label{fig:arch}
\end{figure}

The system uses \code{browser-use} to drive Chromium through CDP. At each
step, the fixed executor receives the task, current DOM-derived page state, and
recent action history, then selects one structured browser action. All runs use
\code{doubao-seed-2-0-pro-260215} through Volcengine Ark, temperature 0.0,
and \code{use\_vision=false}. The navigator is \code{deepseek-chat} at
temperature 0.0 (billed under the provider's DeepSeek-V4-Flash tier).

We compare six configurations. \textbf{E} uses the executor alone. \textbf{I}
injects one opening navigator plan. \textbf{R-$k$} additionally refreshes a
short tactical subgoal every $k\in\{1,3,5\}$ raw agent steps. \textbf{R-A}
re-queries the navigator only after a three-step cooldown when either no new
milestone has occurred for at least three consecutive steps or a previously
visited canonicalized URL is revisited. The navigator is instructed to produce a
short plan and recovery guidance, never browser-action JSON.

\subsection{Tasks, strict success, and run matrix}

Table~\ref{tab:tasks} summarizes four frozen task cards. They span read-only
retrieval, hierarchical issue navigation, and multi-constraint filtering on live,
publicly browsable sites. Each card fixes its prompt, starting conditions,
forbidden actions, recovery rules, expected primary domain, and hard success
criteria.

\begin{table}[t]
\centering
\caption{Evaluation tasks and strict-success requirements. Full task cards are
provided in the supplementary material.}
\label{tab:tasks}
\scriptsize
\setlength{\tabcolsep}{3pt}
\begin{tabularx}{\textwidth}{@{}p{1.55cm}p{1.42cm}p{1.35cm}Y@{}}
\toprule
\textbf{Task} & \textbf{Type} & \textbf{Site} & \textbf{Hard criterion} \\
\midrule
Shopping & Read-only query & \code{amazon.com} &
At least three comparable goods options, each with price, seller, and source URL; no purchase. \\
Hospital & Read-only query & \code{map.baidu.com} &
Three distinct facilities, each with name, phone, address, and source URL; explicitly state missing fields. \\
GitHub & Issue navigation & \code{github.com} &
Filter \code{label:bug}\, \code{is:open}, sort by oldest, open the earliest issue, and report its body and first comment. \\
HuggingFace & Filtered selection & \code{huggingface.co} &
Apply Text Generation, PyTorch, and Chinese filters; sort by downloads; open the top model; report the base model or verified-not-visible. \\
\bottomrule
\end{tabularx}
\end{table}

A run counts as strict success only when every hard criterion is
supported by its final answer, saved evidence, and, where necessary, visited
URLs. The deterministic checker was followed by a frozen criterion-level audit
of saved evidence. For example, the hospital rubric requires all three fields
for three distinct facilities; the GitHub rubric requires oldest-issue and
first-comment evidence; and the HuggingFace rubric verifies the active filters,
download sort, and base-model field. This protocol distinguishes strict success
from an agent merely calling \code{done}.

The run matrix contains six configurations $\times$ four tasks $\times$ ten
repetitions, yielding 240 runs. We treat this matrix as a \emph{case study}:
four live workflows chosen to span read-only retrieval, hierarchical navigation,
and multi-constraint filtering, not as a representative web-agent benchmark.
Runs used headful ephemeral-incognito Chromium
at $1280\times720$, one action per step, \code{max\_steps}=35, and
\code{max\_failures}=3. Runs remained sequential within each task block, but
condition order was randomized rather than following a fixed
E$\rightarrow$I$\rightarrow$R-1$\rightarrow$R-3$\rightarrow$R-5$\rightarrow$R-A
sequence.

\subsection{Diagnostic metrics}

Strict success is the headline outcome. For diagnostic interpretation, we report
median steps, wall-clock duration, stall burden, and state revisit
rate. Each task defines an ordered milestone list $M=(m_1,\ldots,m_K)$; a step
counts as progress when it first satisfies the next unachieved checkpoint.
Stall burden is the fraction of steps before \code{done} that do not
produce a new milestone. State revisit rate is the fraction of steps whose
canonicalized URL (domain and path, without query parameters) has already been
seen. Both metrics are computed from saved \code{history.json} traces and are
descriptive rather than normative scores.

For Hospital lookup, the six checkpoints are: (M1) navigate to a map site;
(M2) submit a hospital search; (M3) reach a results list; (M4) select a
specific facility; (M5) observe a phone number on the page; and (M6) extract
fields or call \code{done}. Each milestone is a deterministic predicate over
the step's canonicalized URL, action type, and page text (implementation in
\code{task\_registry.py}). High stall under E often reflects looping between
M3--M4 listing pages that lack phone details before M5 is reached. For state
revisit rate, URLs are canonicalized to lowercase domain plus path; query
strings and fragments are stripped, so hash-only or parameter-only changes
count as the same state. Milestone lists for the other three tasks appear in
the supplementary material.

% =============================================================================
\section{Results}
% =============================================================================

\subsection{Strict success}

Figure~\ref{fig:strict} summarizes strict-success rates by task and
configuration ($n{=}10$ per cell; Wilson 95\% CIs; $k/n$ annotated per bar).
One-shot planning (I) and adaptive replanning (R-A) each achieved 40/40 strict
successes across the four tasks. The benefits of repeated replanning were not
monotonic: R-1 and R-3 each dropped to 8/10 on HuggingFace, R-1 matched the
E baseline on GitHub (9/10), and R-3 fell to 7/10 on Hospital. Shopping was
at ceiling in every condition and therefore provides little discrimination.
Given $n{=}10$ per cell, intervals are descriptive rather than inferential.

\begin{figure}[t]
\centering
\figplaceholder{figures/fig2.png}{0.94}
\caption{Strict success rate by task and configuration. Bars show Wilson 95\%
CIs; labels give $k/n$. Shopping is at ceiling in all conditions; Hospital and
HuggingFace show the largest configuration spread. Full counts and run-level
logs are provided in the supplementary material.}
\label{fig:strict}
\end{figure}

All 240 runs called \code{done}, yet 17 failed the strict rubric
(E: 8; R-1: 3; R-3: 5; R-5: 1; I: 0; R-A: 0). This discrepancy confirms that
self-report alone is insufficient for evaluating task completion; live-site
variation further limits strong inferential claims.

\subsection{Efficiency and Hospital diagnostic}

Table~\ref{tab:eff} reports median steps, wall-clock duration, estimated
USD cost, and total tokens (executor plus navigator). On Hospital, one-shot
planning reduced median steps from 14 to 7 and duration from 381\,s to 199\,s
while increasing strict success from 3/10 to 10/10; estimated cost fell from
\$0.20 to \$0.10 and tokens from 222k to 114k. R-A retained 10/10 strict
success with the lowest Hospital medians on every efficiency column (5 steps,
145\,s, \$0.08, 85k tokens). Across the full task set, R-A had the lowest
median duration on Shopping, Hospital, and GitHub; R-5 was slightly faster on
HuggingFace (294\,s versus 299\,s for R-A). Navigator cadence also shifts token
and cost overhead: on Shopping, R-1 raised median cost to \$0.22 and tokens to
264k versus \$0.14 and 163k under E, while still reaching 10/10 strict success.
On Shopping, one-shot guidance increased median duration relative to E (339\,s
versus 246\,s) despite comparable token cost (\$0.18 versus \$0.14), suggesting
that navigator-induced exploration can add wall-clock latency even when terminal
outcomes are unchanged. These are run-set summaries, not evidence of universal
per-task optimality.

% Auto-generated by scripts/compute_paper_stats.py

\begin{table}[t]
\centering
\caption{Efficiency metrics by task and navigator cadence (median with 95\%
CI; $n{=}10$ per cell). Token counts in thousands (k). Cost estimates use
provider price schedules at run time. Hospital rows show the largest
cadence-dependent spread; Shopping rows are at strict-success ceiling across
conditions.}
\label{tab:eff}
\footnotesize
\renewcommand{\arraystretch}{1.05}
\setlength{\tabcolsep}{3pt}
\begin{tabularx}{\textwidth}{@{}ll*{4}{>{\raggedleft\arraybackslash}X}@{}}
\toprule
\textbf{Task} & \textbf{Cond.} & \textbf{Steps} & \textbf{Dur.~(s)} & \textbf{Cost (USD)} & \textbf{Tokens} \\
\midrule
\multirow{6}{*}{Shopping} & E & 6~{\scriptsize[6, 7]} & 246~{\scriptsize[242, 329]} & 0.14~{\scriptsize[0.14, 0.17]} & 163k~{\scriptsize[160k, 194k]} \\
 & I & 7~{\scriptsize[6.5, 8.0]} & 339~{\scriptsize[242, 362]} & 0.18~{\scriptsize[0.15, 0.21]} & 208k~{\scriptsize[173k, 240k]} \\
 & R-1 & 8.5~{\scriptsize[7.5, 9.0]} & 438~{\scriptsize[302, 544]} & 0.22~{\scriptsize[0.16, 0.29]} & 264k~{\scriptsize[199k, 333k]} \\
 & R-3 & 9~{\scriptsize[7, 11]} & 297~{\scriptsize[257, 435]} & 0.18~{\scriptsize[0.16, 0.25]} & 206k~{\scriptsize[188k, 285k]} \\
 & R-5 & 8~{\scriptsize[7, 9]} & 288~{\scriptsize[253, 402]} & 0.16~{\scriptsize[0.15, 0.26]} & 192k~{\scriptsize[176k, 295k]} \\
 & R-A & 6~{\scriptsize[6, 7]} & 230~{\scriptsize[211, 253]} & 0.14~{\scriptsize[0.14, 0.16]} & 167k~{\scriptsize[164k, 188k]} \\
\addlinespace
\multirow{6}{*}{Hospital} & E & 14~{\scriptsize[5.5, 16.0]} & 381~{\scriptsize[169, 415]} & 0.20~{\scriptsize[0.08, 0.21]} & 222k~{\scriptsize[82k, 240k]} \\
 & I & 7~{\scriptsize[5, 7]} & 199~{\scriptsize[142, 254]} & 0.10~{\scriptsize[0.07, 0.10]} & 114k~{\scriptsize[79k, 116k]} \\
 & R-1 & 9~{\scriptsize[9, 12]} & 313~{\scriptsize[286, 359]} & 0.13~{\scriptsize[0.13, 0.16]} & 151k~{\scriptsize[150k, 199k]} \\
 & R-3 & 10~{\scriptsize[9, 10]} & 294~{\scriptsize[264, 302]} & 0.13~{\scriptsize[0.12, 0.14]} & 157k~{\scriptsize[140k, 162k]} \\
 & R-5 & 9~{\scriptsize[9, 9]} & 239~{\scriptsize[228, 256]} & 0.12~{\scriptsize[0.12, 0.12]} & 138k~{\scriptsize[137k, 139k]} \\
 & R-A & 5~{\scriptsize[5, 6]} & 145~{\scriptsize[144, 150]} & 0.08~{\scriptsize[0.08, 0.08]} & 85k~{\scriptsize[84k, 96k]} \\
\addlinespace
\multirow{6}{*}{GitHub} & E & 15~{\scriptsize[13.5, 27.0]} & 440~{\scriptsize[364, 687]} & 0.26~{\scriptsize[0.22, 0.36]} & 297k~{\scriptsize[253k, 418k]} \\
 & I & 15.5~{\scriptsize[12, 19]} & 438~{\scriptsize[305, 489]} & 0.26~{\scriptsize[0.21, 0.31]} & 303k~{\scriptsize[237k, 359k]} \\
 & R-1 & 13.5~{\scriptsize[12, 17]} & 463~{\scriptsize[356, 531]} & 0.24~{\scriptsize[0.21, 0.29]} & 290k~{\scriptsize[251k, 347k]} \\
 & R-3 & 11.5~{\scriptsize[11, 13]} & 372~{\scriptsize[371, 431]} & 0.21~{\scriptsize[0.20, 0.23]} & 242k~{\scriptsize[237k, 265k]} \\
 & R-5 & 12~{\scriptsize[11, 13]} & 383~{\scriptsize[283, 407]} & 0.21~{\scriptsize[0.18, 0.23]} & 237k~{\scriptsize[216k, 271k]} \\
 & R-A & 14~{\scriptsize[13, 17]} & 344~{\scriptsize[338, 379]} & 0.23~{\scriptsize[0.22, 0.29]} & 273k~{\scriptsize[259k, 347k]} \\
\addlinespace
\multirow{6}{*}{HuggingFace} & E & 13~{\scriptsize[11.5, 13.0]} & 331~{\scriptsize[259, 339]} & 0.19~{\scriptsize[0.17, 0.19]} & 223k~{\scriptsize[196k, 224k]} \\
 & I & 13.5~{\scriptsize[12, 14]} & 317~{\scriptsize[262, 394]} & 0.22~{\scriptsize[0.18, 0.22]} & 251k~{\scriptsize[214k, 256k]} \\
 & R-1 & 13~{\scriptsize[12.5, 16.0]} & 360~{\scriptsize[317, 406]} & 0.21~{\scriptsize[0.20, 0.25]} & 237k~{\scriptsize[228k, 292k]} \\
 & R-3 & 15~{\scriptsize[13, 20]} & 386~{\scriptsize[259, 419]} & 0.24~{\scriptsize[0.20, 0.32]} & 274k~{\scriptsize[236k, 375k]} \\
 & R-5 & 13~{\scriptsize[11.0, 16.5]} & 294~{\scriptsize[252, 308]} & 0.20~{\scriptsize[0.17, 0.25]} & 234k~{\scriptsize[197k, 296k]} \\
 & R-A & 14~{\scriptsize[12.5, 16.0]} & 299~{\scriptsize[272, 328]} & 0.22~{\scriptsize[0.19, 0.25]} & 253k~{\scriptsize[224k, 291k]} \\
\bottomrule
\end{tabularx}
\end{table}


Hospital provides the clearest process-level diagnostic
(Fig.~\ref{fig:hospital}; Table~\ref{tab:hospital}). The cadence sweep is
non-monotonic: strict success rises from 3/10 under E to 10/10 under I and
R-1, drops to 7/10 under R-3, and recovers to 10/10 under R-A, while median
stall burden moves inversely from 0.71 (E) to 0.17 (R-A). Relative to E, I
sharpens execution---median steps fall from 14 to 7 and duration from 381\,s
to 199\,s---consistent with the opening plan avoiding listing pages that lack
phone details. R-A preserves 10/10 strict success with the lowest steps (5),
duration (145\,s), stall (0.17), and revisit rate (0.33). In contrast, R-1/R-3
raise revisit rates to 0.56--0.60 without improving on I's reliability. This
supports a narrow interpretation: in this recovery-heavy workflow,
event-triggered guidance can preserve one-shot reliability while avoiding some
unproductive navigation. It does not establish that the trigger is optimal or
that the same pattern generalizes to all tasks.

\begin{figure}[t]
\centering
\figplaceholder{figures/fig_hospital_cadence.png}{0.88}
\caption{Hospital lookup: strict success rate and median stall burden across
navigator cadence. Success is non-monotonic in replan frequency (R-3 drops to
7/10); R-A matches I on reliability while lowering stall burden further.
Steps, duration, and revisit medians appear in Table~\ref{tab:hospital}.}
\label{fig:hospital}
\end{figure}

\begin{table}[t]
\centering
\caption{Hospital diagnostic metrics (medians). Stall and revisit are
fractions; strict success is out of 10. Fig.~\ref{fig:hospital} visualizes
strict success and stall across cadence; this table adds steps, duration, and
revisit rate.}
\label{tab:hospital}
\small
\setlength{\tabcolsep}{5pt}
\begin{tabular}{@{}lrrrrr@{}}
\toprule
\textbf{Cond.} & \textbf{Strict} & \textbf{Steps} & \textbf{Dur. (s)} & \textbf{Stall} & \textbf{Revisit} \\
\midrule
E   & 3/10  & 14 & 381 & 0.71 & 0.43 \\
I   & 10/10 & 7  & 199 & 0.29 & 0.43 \\
R-1 & 10/10 & 9  & 313 & 0.44 & 0.56 \\
R-3 & 7/10  & 10 & 294 & 0.50 & 0.60 \\
R-5 & 9/10  & 9  & 239 & 0.44 & 0.56 \\
R-A & 10/10 & 5  & 145 & 0.17 & 0.33 \\
\bottomrule
\end{tabular}
\end{table}

% =============================================================================
\section{Discussion and Limitations}
% =============================================================================

One-shot navigation is the most reliable result of this study. It lifts
Hospital from 3/10 to 10/10 and GitHub from 9/10 to 10/10 (Fig.~\ref{fig:strict})
while maintaining ceiling performance on the other two tasks. Thus, an opening
plan can be useful when the executor must maintain field-completeness or
filter/sort constraints over multiple page transitions.

By contrast, fixed-interval replanning is not uniformly beneficial. Frequent
intervention sometimes disrupts a working search flow: R-1 loses the one-shot
gain on GitHub, while R-1 and R-3 both lose ground on HuggingFace
(Fig.~\ref{fig:strict}). On Hospital, success is non-monotonic in replan
frequency---R-3 falls to 7/10 despite more frequent guidance than R-5
(Fig.~\ref{fig:hospital}). Longer intervals often recover the one-shot baseline,
but the data do not support one globally optimal interval. R-A matches I's
40/40 strict-success record and, in this run set, often lowers median duration
(Table~\ref{tab:eff}), with the strongest diagnostic advantage on Hospital:
10/10 strict success, lowest median steps and duration, and lowest stall and
revisit rates (Fig.~\ref{fig:hospital}; Table~\ref{tab:hospital}). Its role
should therefore be interpreted as a promising recovery policy rather than a
demonstrated universal optimum.

Several limitations bound this interpretation. Although the condition order was
randomized within task blocks, runs remained sequential on live websites, so
content drift, latency, or rate-limiting may still contribute to run-to-run
variation; $n=10$ per cell also limits strong inference. The study evaluates one
text-only deployment with a Doubao executor and a DeepSeek navigator, so it
estimates a system-level intervention rather than a model-independent causal
effect of planning. The four workflows are read-only or filter-based and exclude
authenticated, write, and visually grounded interaction. Finally, R-A uses one
hand-specified stall/revisit trigger without threshold-sensitivity analysis.
Within these bounds, the sweep still isolates navigator cadence under a fixed
executor and evidence-audited scoring; extending the protocol to additional
models, authenticated or visual workflows, and trigger designs remains future
work.

% =============================================================================
\section{Conclusion}
% =============================================================================

Holding a CDP browser agent's executor fixed, this case-study cadence sweep
shows that navigator guidance is neither uniformly helpful nor uniformly harmful.
One-shot planning and event-triggered replanning each achieved 40/40
evidence-audited strict successes in this four-task snapshot, while frequent
fixed-interval replanning reduced success on selected tasks. The Hospital workflow illustrates the clearest benefit
(Fig.~\ref{fig:hospital}; Table~\ref{tab:hospital}): one-shot guidance improved
completion and reduced the median execution burden, and event-triggered
guidance further reduced median steps, duration, stall burden, and revisit rate
in this run set. Future work should introduce same-model
navigator controls, include visual and authenticated workflows, increase the
number and diversity of tasks, and test trigger thresholds more systematically.

\paragraph{Reproducibility.}
All claims trace to frozen task definitions (\code{task\_cards.json}), run logs
(\code{agent\_runs.json}), audit labels (\code{run\_audit.csv}), and the
adjudicator/navigator source (\code{runner.py}, \code{navigator.py},
\code{prompts.py}). Artifacts will be released upon acceptance.

\begin{credits}
\subsubsection{\discintname} The authors declare no competing interests
relevant to this work.
\end{credits}

\bibliographystyle{splncs04}
\bibliography{references}

\end{document}
